\begin{helper}
For every two-bite sample \(\omega\), vertex set \(I\), and tagged projected
coordinate \(e\), same-color closedness in the one-sided
\(C^+(I)\) sense implies full closedness:
\[
\noderef{TwoBiteProjectionPairSameColorClosed}(\omega,I,e)
\Longrightarrow
\noderef{TwoBiteProjectionPairClosed}(\omega,I,e).
\]
\end{helper}
\begin{proof}
We prove the red and blue tag cases separately.

Suppose first that \(e=\operatorname{red}(q)\) and that
\(\noderef{TwoBiteProjectionPairSameColorClosed}(\omega,I,e)\) holds.  By the
definition of \(\noderef{TwoBiteProjectionPairSameColorClosed}\), there is a
red base vertex \(r\) such that
\[
q\in
\noderef{TwoBiteBasePairs}
\bigl(\pi_R^\omega(\noderef{TwoBiteXPlus}(\omega,I,\operatorname{red}(r)))\bigr).
\tag{1}
\]
Unfolding \(\noderef{TwoBiteBasePairs}\), (1) says that both endpoints of the
ordered base pair \(q\) lie in the red projection of
\(\noderef{TwoBiteXPlus}(\omega,I,\operatorname{red}(r))\), with the required
increasing endpoint order.  The endpoint-order part of (1) is unchanged in
the full-closed predicate, so only the endpoint membership has to be enlarged.

Take an arbitrary vertex
\[
v\in\noderef{TwoBiteXPlus}(\omega,I,\operatorname{red}(r)).
\]
By the definition of \(\noderef{TwoBiteXPlus}\), \(v\in I\) and
\[
v\in\noderef{TwoBiteLiftedPositiveNeighborhood}(\omega,\operatorname{red}(r)).
\]
Unfolding \(\noderef{TwoBiteLiftedPositiveNeighborhood}\) in the red case,
this means
\[
r<\pi_R^\omega(v)
\quad\text{and}\quad
\noderef{TwoBiteRedGraph}(\omega)\text{ contains the edge }
\{r,\pi_R^\omega(v)\}.
\]
The second conjunct is exactly the red-case membership condition for
\[
v\in\noderef{TwoBiteLiftedNeighborhood}(\omega,\operatorname{red}(r)).
\]
Together with \(v\in I\), this gives
\[
v\in\noderef{TwoBiteX}(\omega,I,\operatorname{red}(r)).
\]
Therefore
\[
\noderef{TwoBiteXPlus}(\omega,I,\operatorname{red}(r))
\subseteq
\noderef{TwoBiteX}(\omega,I,\operatorname{red}(r)),
\]
and applying the red projection preserves this containment.  Hence each
endpoint of \(q\) which lies in the projected \(X_r^+(I)\) set in (1) also
lies in the projected \(X_r(I)\) set.  Keeping the same endpoint order, (1)
implies
\[
q\in
\noderef{TwoBiteBasePairs}
\bigl(\pi_R^\omega(\noderef{TwoBiteX}(\omega,I,\operatorname{red}(r)))\bigr).
\]
This is precisely the red case of
\(\noderef{TwoBiteProjectionPairClosed}(\omega,I,e)\), witnessed by the same
base vertex \(\operatorname{red}(r)\).

The blue case is identical with the colors interchanged.  If
\(e=\operatorname{blue}(q)\), the one-sided same-color closed predicate gives
a blue vertex \(b\) and membership of \(q\) in the blue projected pairs of
\(\noderef{TwoBiteXPlus}(\omega,I,\operatorname{blue}(b))\).  Membership in
the positive blue lifted neighborhood consists of the larger-index condition
\(b<\pi_B^\omega(v)\) together with the blue base adjacency condition.  Dropping
the larger-index condition leaves exactly membership in the full blue lifted
neighborhood, so
\[
\noderef{TwoBiteXPlus}(\omega,I,\operatorname{blue}(b))
\subseteq
\noderef{TwoBiteX}(\omega,I,\operatorname{blue}(b)).
\]
The same projection and endpoint-order argument shows that \(q\) is a blue
projected pair inside \(X_b(I)\).  Thus
\(\noderef{TwoBiteProjectionPairClosed}(\omega,I,e)\) holds in the blue case
as well.
\end{proof}
